\begingroup
\selectlanguage{english}
\noindent\textbf{Cofinality, connected objects, and the recovery of diagrams.}

Chapter~II begins by asking when a change of indexing category leaves sections unchanged \citep{Grothendieck1990II}. It ends this first part with the stronger question of when restriction recovers a whole diagram and every natural transformation between diagrams. Between these questions Grothendieck develops a topos-theoretic account of filtering, examines the homotopy information lost by ordinary connectedness, and specializes the language of derivators to sets. The movement is deliberate. Ordinary cofinality is not merely an elementary preliminary to the derived theory: its proof identifies the object which must be derived, and its limitations identify the higher information which a derivator must retain.

The conventions are essential. A coefficient on $C$ is a contravariant functor, and the category denoted $c\backslash C'$ has objects $c\to\alpha(c')$. Thus a cofinal functor in these pages is final in the usual covariant terminology; its opposite is initial for the diagrams whose limits are being taken. Under this convention the first question has a representability answer. The presheaf
\[
 K_\alpha(c)=\pi_0(c\backslash C')= (\alpha_!1_{C'})(c)
\]
represents $F\mapsto\Gamma(C',\alpha^*F)$. The comparison with $\Gamma(C,F)$ is induced by $K_\alpha\to1_C$. It is an isomorphism for all $F$ exactly when every indicated comma category is nonempty and connected. This is the classical cofinality theorem of SGA~4, Expos\'e~I, \S8.1 \citep{SGA4I1972}. The initial weak cofiltering hypothesis is sufficient, and Grothendieck replaces it by the exact connectedness criterion as the argument develops.

The same proof suggests its own derived form. For a commutative ring $k$, replace the component presheaf by the derived weight $\mathbf L\alpha_!k_{C'}$. Its value at $c$ is represented by chains on $N(c\backslash C')$. Derived adjunction then gives
\[
 \mathbf R\Gamma(C',\alpha^*F)
 \simeq\mathbf R\operatorname{Hom}_{C;k}
       (\mathbf L\alpha_!k_{C'},F).
\]
Consequently weak contractibility of all the comma nerves implies cohomological cofinality, answering the question on manuscript page~5 affirmatively. For these fixed linear coefficients, their $k$-acyclicity already suffices. For spaces the weight is the corresponding presheaf of classifying spaces, and the mapping-space adjunction gives the homotopical statement. Modern higher category theory places this argument within the general criterion for cofinality \citep[Theorem~4.1.3.1]{Lurie2009HTT}. The commentary constructs the linear comparison by a bar resolution, so the answer does not depend on the internal-Hom formula subsequently proposed in the manuscript.

That formula is a genuine point of failure. The global derived Hom belongs to the category of coefficients at a point; a sheaf-valued internal Hom belongs to the diagram category. Applying global sections can identify two such objects without identifying the objects themselves. For the terminal-vertex inclusion $i:\{1\}\hookrightarrow[1]$, a presheaf is an arrow $F_1\to F_0$, and
\[
 i_!A=(A\xrightarrow{1}A),\qquad i_*A=(A\to *).
\]
The comma categories for cofinality are points, and global sections are $F_1$. Nevertheless $i_*i^*F$ need not be $F$, while $\underline{\operatorname{Hom}}(i_!1,F)=F$. This disproves the unrestricted internal identification in sets, hence already before derived complications arise. The correct internal comparison is equivalent to an additional projection-formula property, whose canonical arrow has the opposite direction from the one displayed in the manuscript. Our point-inclusion calculation identifies its obstruction exactly: for a fixed presheaf, every restriction along an arrow into the chosen object must be invertible. The formula holds for all coefficients precisely when all those incoming arrows are isomorphisms.

The turn to toposes is motivated by a different problem: how to recognize the equalization part of filtering without inspecting individual arrows. An irreducible presheaf topos records common predecessors of objects, a condition visible in the underlying preorder. Local irreducibility records common predecessors compatible with a prescribed pair of arrows, namely PS~1. It still does not supply PS~2, the equalization of arbitrary parallel arrows. The parallel-arrow category $0\rightrightarrows1$ separates the first two conditions. A nontrivial group category $BG$ separates local irreducibility from local cofilteredness: every representable slice is equivalent to a point, but cancellation prevents the equalization of distinct group elements. Several preliminary identifications in the manuscript are explicitly withdrawn by its annotations. The audit preserves this sequence of correction and distinguishes it from errors of transcription.

Total $0$-connectedness adds preservation of products by the component functor. For $\widehat X=\operatorname{Fun}(X^{\mathrm{op}},\mathrm{Set})$, it says that $X^{\mathrm{op}}$ is an ordinary sifted category. Sifted colimits preserve finite products; filtered colimits preserve finite limits, including equalizers \citep{AdamekRosickyVitale2010}. The simplex category supplies the instructive distinction. The products of its representables have connected categories of elements, but the two vertices of $[1]$ admit no common predecessor over $[1]$. Thus its presheaf topos is totally $0$-connected without being locally irreducible. Nothing about this ordinary product test alone makes all the relevant nerves weakly contractible.

The surviving theorem in 2.6 makes the product condition hereditary under connected slices. Its equivalent formulations are (a), (a'), (b), and (d); the generator-slice conditions (c) and (c') are strictly weaker. The diagonal $b\to b\times b$ explains the distinction. Once $b\times b$ is connected and the topos is locally irreducible, its pullback against $(u,v):a\to b\times b$ is noninitial and yields an equalizing arrow. Conversely, PS~1 and PS~2 lift to categories of elements, giving connected products in every connected slice. The commentary supplies the omitted generator-component lemma and the successive equalizations in this argument. It also restores connectedness of the base where the printed formulation of 2.8 omits it.

On a connected topos satisfying this hereditary condition, $\pi_0$ preserves finite limits. This makes the later cohomological argument precise. The induction must run over connected, locally totally $0$-connected toposes, since it repeatedly passes to connected components of slice toposes. The degree-zero row of the covering spectral sequence is the augmented complex of functions $I^{p+1}\to M$, where $I$ is the component set of a covering object. Inserting a chosen element of $I$ contracts that complex. The stronger homotopical conclusion follows from hypercovers: applying the exact component functor to their matching-object diagrams gives boundary fillers in every degree, so the resulting simplicial sets are weakly contractible. This connects the argument to the classical hypercover construction of homotopy types \citep[\S\S8--9]{ArtinMazur1969}. It concerns constant coefficients and shape. It does not trivialize arbitrary sheaf cohomology, as the worked calculation $H^1(\mathbf P^1_k,\mathcal O(-2))\cong k$ illustrates.

The ind-object reformulation in 2.15--2.18 returns to a classical theorem: a presheaf is a filtered colimit of representables exactly when its category of elements is filtered \citep[Expos\'e~I, Theorem~8.3.3]{SGA4I1972}. Taking the terminal presheaf recovers filteredness of the indexing category. Taking its ordinary left Kan extension along an arbitrary functor shows that the resulting category of elements is again filtered. The modern minimal-localizer theorem supplies the asphericity conclusion for every fundamental localizer, without the additional limit axiom used in the historical formulation \citep[Theorem~2.2.11]{Cisinski2004Minimal}.

The converse question on page~36 has a stronger quantifier than later uses of the phrase ``totally aspherical''. In the commentary we denote by $T(A)$ the condition that every ordinary cofinal target of $A$ have weakly contractible nerve. Filtered categories satisfy it. The attempted proof that ordinary siftedness suffices is retracted on pages~36bis--36ter; our explicit completion of the counterexample uses the simplicial circle and a one-point collage whose classifying space is $S^2$. We also prove a necessary common-successor condition for $T(A)$. Its obstruction is a four-object category with a cofinal three-object subcategory and a nonzero integral two-cycle. This settles $T(A)\Leftrightarrow A\text{ filtered}$ for preorders. For general categories the remaining issue is equalization of parallel arrows. We did not locate a resolution under this exact universal quantifier in the sources checked; this is a report of that bounded search, not a declaration of a recognized open problem. Results developed here are marked as the commentators' work and carry no claim of priority.

Section~3 shows how much of the earlier derivator terminology can be tested with sets. Pointwise localization of category-valued diagrams at component equivalences yields the represented derivator of sets. Properness and smoothness become connectedness of strict lift categories. A universal $W_0$-equivalence can be tested over vertices and arrows; a strong $W_0$-fibration is characterized component by component by those universal equivalences. Coverings explain why the weak condition is weaker: monodromy can connect different fibre components without making each fibre connected. The comparison with the commentary on Chapter~I is made at the precise lifting and monodromy arguments, rather than repeating its whole discussion.

The final transition is from preservation of global sections to full faithfulness of restriction. Proposition~4.1 characterizes the latter by nonempty connected categories of factorizations
\[
 b_0\longrightarrow u(a)\longrightarrow b_1
\]
of every arrow $b_0\to b_1$. These are the lax epimorphisms, or absolutely dense functors, of the later literature \citep{AdamekElBashirSobralVelebil2001}. The homotopical version requires weakly contractible factorization categories \citep[Theorem~6.10]{GrothStovicek2016Trees}. A quotient of groups gives a particularly concrete distinction: the factorization category has the homotopy type of the classifying space of the kernel. Two qualifications belong to the audit. The printed coslice condition in 4.1(d') must be initial, not final. Also the coefficient-specific discussion on page~51 invokes pointwise Kan-extension formulas without the necessary comma limits; a discrete coefficient category on two objects gives a counterexample even when the indexing category is ordered. The corrected theorem quantified over all coefficient categories is established by its set-valued case and Yoneda embeddings.

The historical setting is the program of coherent diagram categories described in \citet[\S69]{Grothendieck1983}, alongside the systematic homotopy Kan extensions of \citet{Heller1988}. The commentary keeps the successive formulations, counterexamples, and marginal revisions visible. It distinguishes a proof supplied here from a proof actually written in these pages, and distinguishes a settled theorem from an unsuccessful search for a later treatment.

\medskip
\noindent\small\textit{AI-produced commentary: GPT-6 Astra (OpenAI, ChatGPT Pro mode). The project's supplied editorial credit is ``edited with Claude (Anthropic)''; no additional Claude editing was performed during this session.}
\endgroup
